\documentclass[../Main.tex]{subfiles}

\begin{document}

\section{Schlussfolgerungen und weiteres Vorgehen} \label{sec:Schlussfolgerungen}

Damit sind wir am Schluss unserer Arbeit angekommen. In diesem Kapitel versuchen wir, unsere Erkenntnisfrage zu beantworten, und gehen dabei auf die Limitationen ein, die wir hatten. Zudem liefern wir weitere Ansätze zur Beantwortung unserer Erkenntnisfrage und schlagen Verbesserungen für unseren Prototyp vor.

\subsection{Beantwortung der Erkenntnisfrage} \label{sec:Beantwortung}

Unsere Erkenntnisfrage lautete:

\begin{quote}
    \textit{Wie müssen Gegnerverhalten und Gegnerplatzierung in einem 3D-Hack-and-Slash-Prototyp konzipiert sein, damit ein zusammenhängender Kampfabschnitt mit nachvollziehbarem Spielfluss, angemessener Schwierigkeit und befriedigendem Spielgefühl entsteht?}
\end{quote}

Im Folgenden fassen wir zusammen, was unsere acht Designvorschläge zur Beantwortung beitragen. \textcolor{red}{Wir ordnen sie dabei danach, wie deutlich die Beobachtungen aus Abschnitt~\ref{sec:ErgVorschlag} für ihre beabsichtigte Wirkung sprechen.}

Für das Gegnerverhalten zeigt sich bei der Begrenzung der gleichzeitig erscheinenden Gegner (V2) und bei der Kombination von Nah- und Fernkämpfern (V3) eine starke Tendenz zur beabsichtigten Wirkung. Dadurch werden nämlich die grossen Gegnermengen, die es in 3D-Hack-and-Slash-Spielen gibt, bekämpfbar, und die Spieler müssen entscheiden, welchen Gegner sie zuerst angreifen. \textcolor{red}{Ebenso deutlich sprechen unsere Beobachtungen dafür, dass sich der Spieler als Grundzustand überlegen fühlt (V1) und dass das dynamische Verhalten der Gegner Druck erzeugt (V5).}

\textcolor{red}{Bei der Gegnerplatzierung zeigt sich eine starke Tendenz dafür, dass die gestaffelte Einführung neuer Gegner wie beabsichtigt wirkt (V8):} Die neuen Gegner wurden erkannt, und ihre spätere Kombination hat die Schwierigkeit erhöht. \textcolor{red}{Gleiches gilt für die Pausen zwischen den Räumen (V7) und für die weiten Räume, die mehr Bewegungsfreiheit und taktische Distanz bieten (V6).}

\textcolor{red}{Für das gezielte Durchbrechen der Überlegenheit (V1) zeigt sich dagegen eine erkennbare, aber keine starke Tendenz zur beabsichtigten Wirkung.} Das liegt daran, dass das Durchbrechen nicht immer beabsichtigt war und nicht immer Spannung, sondern auch Frust erzeugt hat. Dasselbe gilt dafür, dass die Schwierigkeit taktisch statt mechanisch entsteht (V3), für die lesbare Anfangsaufstellung (V5), für die engen Räume (V6), für die Pausen zwischen den Wellen und den planbaren Anstieg der Schwierigkeit (V7) sowie dafür, dass die isolierte Einführung beim Verstehen neuer Gegner hilft (V8).

Bei der Ankündigung von Angriffen (V4) deutet sich an, dass bei den Nahkämpfern Animationen allein ausreichen, bei den Archers finden wir dagegen kaum Hinweise darauf. Ob eine Ankündigung über weitere Kanäle wie Sound bei den Archers ausreicht, konnten wir nicht überprüfen.

Durch all diese Designvorschläge zeigt sich eine starke Tendenz dafür, dass unsere Schwierigkeit angemessen war. \textcolor{red}{Auch für einen gelungenen Spielfluss sprechen unsere Daten deutlich:} Vier von fünf Testpersonen erlebten den Kampfabschnitt als zusammenhängend, und zwei spielten nach dem Ende sogar ungefragt weiter. \textcolor{red}{Für ein befriedigendes Spielgefühl zeigt sich dagegen eine erkennbare, aber keine starke Tendenz, da} es durch Faktoren beeinflusst wurde, die nicht direkt mit den Gegnern zu tun hatten, wie zum Beispiel, dass die Combos etwas zu träge waren (Abschnitt~\ref{sec:Gesamteindruck}). Auch der Spielfluss wurde an einer Stelle durch einen solchen Faktor gebrochen, weil das Ende eines Encounters nicht kommuniziert wurde.

Deshalb lautet die Antwort auf unsere Erkenntnisfrage vorläufig: Für unseren Prototyp deuten die Playtests darauf hin, dass Gegnerverhalten und Gegnerplatzierung vor allem dann zu einem nachvollziehbaren Spielfluss, einer angemessenen Schwierigkeit und einem befriedigenden Spielgefühl beitragen, wenn die Zahl der gleichzeitig aktiven Gegner begrenzt ist, Nah- und Fernkämpfer kombiniert werden, neue Gegner isoliert eingeführt und erst danach kombiniert werden und es zwischen den Räumen Pausen gibt. Die lesbare Anfangsaufstellung und die Pausen zwischen den Wellen konnten wir dagegen nur in Ansätzen belegen.

Diese Arbeit entstand aus dem Wunsch, Spiele nicht nur zu spielen, sondern sie genauer zu verstehen und selbst ein solches Spielerlebnis zu erschaffen. Das haben wir mit diesem Prototyp erreicht, und zumindest für unseren Kampfabschnitt konnten wir Tendenzen erkennen, wie Gegnerverhalten und Gegnerplatzierung zu einem solchen Spielerlebnis beitragen. Zudem betreffen diese Tendenzen speziell das 3D-Hack-and-Slash-Genre.

\subsection{Limitationen} \label{sec:Limitationen}

\textcolor{red}{Die Ergebnisse dieser Arbeit unterliegen mehreren Einschränkungen.}

\textcolor{red}{Bei den Playtests ist zunächst die \textcolor{red}{kleine Stichprobe [Begriff: \enquote{Stichprobe} meint in der Statistik eine nach Regeln gezogene Auswahl. Eure Testpersonen habt ihr gezielt ausgesucht. Vorschlag: \enquote{die geringe Zahl der Testpersonen}]} zu nennen: Mit fünf Testpersonen lassen sich nur Tendenzen erkennen (Abschnitt~\ref{sec:Auswertung}). Die Urteile in Kapitel~\ref{sec:Evaluation} \textcolor{red}{sind deshalb als Tendenzen und nicht als Bestätigungen zu verstehen. Wo wir kaum oder keine Hinweise gefunden haben,} bedeutet das nicht, dass das Gegenteil zutrifft. Zudem kannten wir die meisten Testpersonen persönlich oder über Bekannte, was trotz unserer Bitte um ehrliche Antworten zu wohlwollenderen Einschätzungen geführt haben kann (Abschnitt~\ref{sec:AuswahlTP}). Zwei Playtests fanden über Discord statt, wo wir nur Bildschirm und Stimme der Testpersonen wahrnehmen konnten, nicht aber Mimik und Körpersprache (Abschnitt~\ref{sec:Beobachtungsraster}).}

\textcolor{red}{Auch das Post-Questionnaire hat Schwächen. Q5 enthält zwei Aussagen auf einmal (überschaubare Anfangsposition und danach unvorhersehbareres Verhalten), sodass unklar bleibt, auf welchen Teil sich eine Bewertung bezieht.}

\textcolor{red}{Bei der Referenzspielanalyse wurde jedes Spiel nur von einer Person analysiert. Die Struktur des MDA-Fragebogens ist zwar einheitlich, Einschätzungen wie \enquote{befriedigend} bleiben aber subjektiv und personenabhängig (Abschnitt~\ref{sec:Analysefragebogen}). Zudem konnten wir die Funktionsweise der Referenzspiele nur aus dem beobachteten Verhalten \textcolor{red}{erschliessen}, da wir keinen Einblick in deren Quellcode hatten.}

\begin{itemize}[leftmargin=1.5cm]
    \item \textcolor{red}{Diesen Abschnitt habe ich aus euren Angaben geschrieben. Falls gewünscht noch ergänzen: mündliche Erklärung der Steuerung vor dem Spielen (statt Selbstentdecken), Gespräche in entspannten Phasen, Leveldesign-Quellen nicht genrespezifisch, Zeitmangel bei der Auswahl der Recherchepunkte.}
\end{itemize}

\begin{itemize}[leftmargin=1.5cm]
    \item \textcolor{red}{[Notiz: offene Punkte aus Kapitel~\ref{sec:Evaluation} in drei Gruppen thematisieren.] (1) Fragen, die nicht genau messen, was der Vorschlag will: Q5 fragt Anfangsaufstellung und Verhalten gleichzeitig ab (V5). Q6 fragt nur, ob sich enge und offene Räume anders anfühlten, nicht, ob sie schwerer waren (V6). Q8 fragt nur, ob die Staffelung bemerkt wurde, nicht, ob sie beim Verstehen half (V8). Q10 fragt nach einem zusammenhängenden Spielfluss, nicht danach, ob man weiterspielen wollte (Spielfluss). (2) Vorschläge, die nur teilweise umgesetzt wurden und deshalb wenig zeigen konnten: Die Begrenzung der gleichzeitig angreifenden Gegner haben wir kaum genutzt (V2), und Angriffe werden nur über Animationen angekündigt, ohne Sound oder Partikeleffekte (V4). (3) Grenzen der Methode: Mit fünf Testpersonen und ohne Vergleichsversion lassen sich Effekte, die sich überlagern, kaum trennen, etwa ob der Druck vom dynamischen Verhalten (V5) oder vom gleichzeitigen Angreifen (V2) kam. Im Ausblick (Abschnitt~\ref{sec:Verbesserungen}) einen Satz ergänzen, dass diese Fragen beim nächsten Playtest genauer formuliert würden. Optional in Kapitel~8 als Lehre: Fragebögen vorher an einer Person ausprobieren.}
\end{itemize}

\subsection{Weiteres Vorgehen und Ausblick} \label{sec:Verbesserungen}

In einem weiteren Vorgehen würden wir weiterhin Gegnerverhalten und Gegnerplatzierung untersuchen. Zum einen würden wir einige unserer Designvorschläge, bei denen wir \textcolor{red}{nur eine erkennbare Tendenz oder kaum Hinweise} feststellen konnten, noch einmal und spezifischer überprüfen, damit wir diesmal klarere Antworten erhalten. \textcolor{red}{Zum anderen würden wir einzelne Vorschläge abwandeln. Für jeden Designvorschlag sieht das wie folgt aus.}

\paragraph{\textcolor{red}{V1: Überlegenheit als Grundzustand}}
Die Überlegenheit des Spielers konnten wir der Tendenz nach schon recht gut etablieren, sie aber noch nicht richtig geplant brechen \textcolor{red}{(Abschnitt~\ref{sec:ErgVorschlag}, V1)}. Die konkrete Verbesserung wäre dabei, den Heavy Soldier wieder etwas stärker zu machen. \textcolor{red}{Wir würden den Vorschlag weiterverfolgen, da unsere Daten zumindest in diese Richtung tendieren.}

\begin{itemize}[leftmargin=1.5cm]
    \item \textcolor{red}{[Aus V1 ergänzen:] Untersuchen, welche Arten von Brüchen der Überlegenheit Spass machen und welche nicht. Beispiel: Das Heilen hat als taktischer Bruch Spannung erzeugt, ob es Spass gemacht hat, ist offen. Die Trägheit der Combos war dagegen kein taktischer Bruch und eher frustrierend (Abschnitt~\ref{sec:ErgVorschlag}, V1).}
\end{itemize}

\paragraph{\textcolor{red}{V2: Begrenzung gleichzeitig aktiver Gegner}}
\textcolor{red}{Die Begrenzung der gleichzeitig angreifenden Gegner haben wir kaum genutzt. In einem nächsten Playtest würden wir sie gezielt einsetzen, auch um dem Stunlock entgegenzuwirken, der gerade dann entstand, wenn alle Gegner gleichzeitig angriffen (Abschnitt~\ref{sec:ErgVorschlag}, V2).}

\paragraph{\textcolor{red}{V3: Ein taktisches Dilemma durch Nah- und Fernkämpfer}}
\textcolor{red}{Das Dilemma wurde wahrgenommen, liess sich aber auch ignorieren (Abschnitt~\ref{sec:ErgVorschlag}, V3).}

\begin{itemize}[leftmargin=1.5cm]
    \item \textcolor{red}{[Fehlt:] Was würdet ihr bei V3 ändern oder genauer prüfen? Zum Beispiel: ob das Dilemma stärker wirkt, wenn Archers mehr Schaden machen und man sie deshalb nicht mehr ignorieren kann, oder eine Frage, die direkt nach taktischer statt mechanischer Schwierigkeit fragt.}
\end{itemize}

\paragraph{\textcolor{red}{V4: Klar telegrafierte Angriffe}}
Die Ankündigung von Angriffen würden wir mit mehreren Kanälen erneut testen\textcolor{red}{, etwa mit Sound und Partikeleffekten}. Dabei könnten wir vielleicht auch Unterschiede bei den Soldiers feststellen, ganz besonders aber bei den Archers. \textcolor{red}{Auch diesen Vorschlag würden wir weiterverfolgen, da unsere Daten zumindest bei den Nahkämpfern in diese Richtung tendieren.}

\paragraph{\textcolor{red}{V5: Lesbare Platzierung, dynamisches Verhalten}}
\textcolor{red}{In einem nächsten Schritt würden wir das dynamische Verhalten weglassen und nur die symmetrische Platzierung überprüfen.} Das ist uns dieses Mal nicht ganz gelungen, weil dieser Vorschlag zwei unterschiedliche Dinge zu überprüfen versuchte. Zudem könnte man prüfen, ob es einen Unterschied macht, ob die Gegner schon da sind, wenn man den Raum betritt, oder ob sie erst dann spawnen\textcolor{red}{, wie in \game{NieR: Automata} bzw.\ \game{Ninja Gaiden} (Abschnitt~\ref{sec:ErgVorschlag}, V5)}.

\paragraph{\textcolor{red}{V6: Der Raum als Druckmittel}}
\textcolor{red}{Die weiten Räume haben wie beabsichtigt gewirkt. Bei den engen Räumen ist die Wirkung dagegen weniger klar und hängt stark vom Spielstil ab. Den Vorschlag würden wir deshalb abgewandelt weiterverfolgen: Da der Druck in engen Räumen vor allem durch den fehlenden Platz zum Heilen entstand (Abschnitt~\ref{sec:ErgVorschlag}, V6), würden wir die Enge gezielter über die Zahl der Gegner steuern und nicht nur über die Raumgrösse. Zudem würden wir die Schwierigkeit der Räume direkt abfragen und nicht nur, ob sie sich anders anfühlten.}

\paragraph{\textcolor{red}{V7: Wellenstruktur mit Erholungsphasen}}
Auch die Wellenstruktur würden wir weiterverfolgen und dabei noch gezielter auf Ruhe und Anspannung achten, zum Beispiel mit einer Pulsmessung während des Playtests.\autocite[285--287]{drachen2018gamesuserresearch} Dabei würden wir zwischen Pausen unterscheiden, die wir dem Spieler vorgeben, und solchen, die sich der Spieler selbst verschafft. \textcolor{red}{Damit die Pausen zwischen den Wellen besser erkennbar sind, bräuchte es zudem eine Rückmeldung, wann eine Welle oder ein Encounter vorbei ist (siehe unten).}

\paragraph{\textcolor{red}{V8: Gestaffelte Einführung neuer Gegner}}
\textcolor{red}{Zwischen der Einführung eines neuen Gegners und seiner Kombination mit anderen würden wir mehr Zeit lassen, etwa wie in \game{NieR: Automata}, wo jede Vorstellungsphase ein ruhiges Tal vor der nächsten Arena bildet (Abschnitt~\ref{sec:ErgVorschlag}, V8). Zudem würden wir genauer abfragen, ob die isolierte Einführung beim Verstehen hilft, und nicht nur, ob sie bemerkt wurde.}

\paragraph{\textcolor{red}{Playtests}}
\textcolor{red}{Allgemein würden wir die Aussagen Q5, Q6, Q8 und Q10 beim nächsten Playtest genauer formulieren, da sie nicht immer genau das abfragten, was der jeweilige Vorschlag bewirken sollte (Abschnitt~\ref{sec:Limitationen}).} Ausserdem könnten wir Playtests mit einer grösseren Gruppe durchführen, um statistisch auswertbare Daten zu erhalten.

\paragraph{\textcolor{red}{Verbesserungen am Prototyp}}
Um dieses Vorgehen zu unterstützen \textcolor{red}{und Faktoren auszuschliessen, die nicht direkt mit den Gegnern zu tun haben}, würden wir an unserem Prototyp folgende Verbesserungen vornehmen:

\begin{itemize}[leftmargin=1.5cm]
    \item die Trägheit der Combos verringern (V1 und Abschnitt~\ref{sec:Gesamteindruck}),
    \item den Übergang zwischen Bewegung und Kampf verbessern (Abschnitt~\ref{sec:Gesamteindruck}),
    \item Audio- und/oder Partikeleffekte ergänzen, die Angriffe ankündigen (V4),
    \item einen Indikator einbauen, der anzeigt, wann ein Encounter vorbei ist, sowie Sounds für die Türen, damit auch diese klar erkennbar sind (Abschnitt~\ref{sec:Gesamteindruck}),
    \item I-Frames hinzufügen, etwa beim Aufstehen nach einem Treffer, um dem Stunlock entgegenzuwirken (V1 und V2),
    \item den Air Combat, der gegen unsere Bodengegner kaum nützlich war und trotzdem von einigen Testpersonen eingesetzt wurde (TP2, TP3 und TP4, jeweils B13), mit zwei Moves ergänzen: einer Aktion, die Gegner in die Luft wirft, und einer, mit der man aus der Luft auf den Boden dashen kann. Einen solchen Angriff aus der Luft wünschten sich auch mehrere Testpersonen (TP1, Q13; TP2, Q13; TP4, B11).
\end{itemize}

\end{document}
